%!TEX root = ../csuthesis_main.tex

\section{总结与展望}
\label{chap:conclusion}

本章对全文研究工作加以总结，梳理本文在基于时序编码的量子时频分析神经网络方向上的主要成果，同时指出当前研究中尚存的不足，并结合量子人工智能的发展态势对后续研究方向作出展望。

\subsection{本文总结}
\label{sec:summary}

本文针对现有经典深度学习音频分类模型参数量庞大，以及现有量子音频处理方法往往忽略信号时序演化特性且在NISQ设备上深层线路训练困难的核心问题，开展了基于量子时序编码的时频分析神经网络构建，以及面向该网络的“经典-量子”跨计算范式训练框架的系统性研究。

第一章首先阐述了研究背景与意义，指出了音频分类在边缘计算等资源受限场景下面临的轻量化需求，分析了当前主流静态量子编码和变分量子线路在处理非平稳时序信号时的局限性，并明确了本文的研究目标与技术路线，旨在构建一套高效表征时序特征并具备良好收敛性的量子神经网络架构与优化框架，以填补量子机器学习在复杂动态序列处理领域的空白。

第二章深入剖析了支撑本研究的相关理论基础与关键技术。本章详细解析了量子计算的基础物理机制，包括量子比特、叠加纠缠特性以及测量原理，明确了其在高维希尔伯特空间中并行处理信息的底层优势；深入阐释了常见量子特征映射理论，重点对比了振幅编码与角度编码的优劣，并剖析了NEQR将经典信息映射为量子基态的内在机理；详细介绍了量子神经网络的端到端架构体系与参数化量子线路的运行机制，并全面梳理了知识蒸馏中“教师-学生”架构、软标签及温度参数调控理论，为后续跨计算范式的模型设计与算法优化构建了坚实的理论支撑。

第三章详细设计了量子时频分析神经网络，构建了端到端量子信息处理流水线。在编码阶段，提出基于双量子比特纠缠的量子时序编码方法，实现了MFCC与时间索引的全局映射，解决了传统编码难以保留动态时序特征的难题。在网络架构上，设计了包含量子卷积、池化与全连接层的演化线路，利用参数化量子门重构局部感受野与降维机制。实验表明，该模型在GTZAN数据集上准确率达75\%，较角度与振幅编码分别提升8\%与12\%；在同等极低参数量下，表现优于QCNN约3\%、高于经典CNN约5\%。此外，结构消融实验证实，若将参数化量子门替换为常规CNOT门，模型准确率将显著下降7\%并陷入局部最优，反向印证了具备连续旋转自由度的参数化线路在捕获非平稳音频特征方面具有不可替代的优势。

第四章系统性地验证了基于第三章提出的量子时频分析神经网络的跨计算范式知识蒸馏优化策略。针对NISQ设备上变分量子线路易陷入局部最优与梯度消失的“贫瘠高原”难题，本章创新性地构建了“经典CNN教师-量子学生”异构蒸馏框架。该策略以在梅尔频谱特征上充分收敛的经典卷积神经网络作为教师端，通过引入温度参数提取包含丰富类别间相似性信息的软标签（暗知识），并结合硬交叉熵与软KL散度构建联合损失函数，对仅含33个可训练参数的量子学生网络进行反向引导。在不增加量子线路参数的前提下，该策略将模型准确率进一步提升至76.17\%，召回率达82.67\%，AUC高达0.8319，显著提升了纯量子模型的收敛速度及抗噪鲁棒性。

综上所述，本文完整构建了一套面向音频分类的量子时频分析神经网络及跨范式训练框架，有效解决了量子音频处理中时序特征流失、参数利用率低及深层线路寻优困难三大难题。研究成果不仅提升了量子机器学习在复杂非平稳序列建模上的表征能力与泛化水平，也为资源受限边缘设备上的轻量化智能音频感知系统提供了底层算法支撑与理论参考。

\subsection{未来工作展望}
\label{sec:future_work}

尽管本文在基于量子计算的音频分类领域进行了系统性探索并取得了突破性的实验验证，但受限于当前量子软硬件的整体发展水平及经典模拟器的算力瓶颈，本研究仍存在部分局限性。未来的研究工作可从以下两个维度进一步拓展：

（1）突破底层量子硬件的规模与物理噪声限制，提升模型在真实设备上的工程可用性。
受限于当前 8 量子比特的硬件设定与理想模拟框架约束，本文在特征提取阶段对高维特征进行了较为激进的降维离散化处理，且暂未深入考量真实的退相干与门误差干扰。未来随着真实物理比特数量的增加及张量网络模拟技术的突破\citess{arute2019quantum}，下一步工作需将底层硬件的拓扑结构与噪声模型纳入考量，引入零噪声外推或概率误差消除\citess{temme2017error}等先进的误差缓解技术。在此基础上，探索将连续的高维梅尔声谱特征直接映射入量子态，在真实超导或离子阱设备上实现针对大规模复杂声学数据集\citess{purwins2019deep}的端到端量子声学建模与多分类验证。

（2）探索量子时序特征提取与自注意力机制的深度融合。
目前的量子时频分析神经网络架构主要依赖量子卷积操作与池化层进行时频特征的局部聚合。经典深度学习在音频序列建模中已证明，全局自注意力机制\citess{vaswani2017attention}在捕获长程依赖方面具有不可替代的优势。未来将尝试把量子自注意力模块嵌入到当前的演化网络中，通过线路演化动态赋予不同时间帧量子态不同的权重分布，从而进一步逼近甚至超越目前最先进的经典声学大模型。